\scanpage{39}
\uline{Formes de} $\mathbb D_n$

Soit $\Gamma$ un groupe, \deleted{tel que} satisfaisant les conditions
suivantes : il existe un ss-groupe $\Gamma^+$ de $\Gamma$ tel que

a) $\Gamma^+$ d'indice 2 dans $\Gamma$, i.e. $\Gamma$ ext. de $\mathbb Z/2\simeq\Gamma/\Gamma^+$
par $\Gamma^+$

b) $\Gamma^+$ \deleted{\unsure{par conséquent}} comm. [dc $\Gamma/\Gamma^+=\mathbb Z/2$ \uline{opère} sur $\Gamma^+$]

c) $\mathbb Z/2$ opère sur $\Gamma^+$ par $x\mapsto-x$

\inserted{d) $\exists\ \sigma\in{}_2\Gamma,\ \sigma\notin\Gamma^+$ [i.e l'extension de $\mathbb Z/2$ par $\Gamma^+$ est scindée].}

e) $\Gamma^+\setminus{}_2(\Gamma^+)$ engendre le groupe $\Gamma^+$. [\uline{NB} Si $\Gamma^+$ est
cyclique, \deleted{ce} d'ordre $n$, ceci signifie $n\ne1$]

\deleted{On a alors}

\deleted{Alors le} Soit Alors $\Gamma^-=\Gamma\setminus\Gamma^+$. On voit tout de
suite \deleted{que} (utilisant a) à d))

\[
\Gamma^-\subset{}_2\Gamma
\]

donc

\[
{}_2\Gamma=\Gamma^-\amalg{}_2\Gamma^+
\]

donc

\[
\Gamma\setminus{}_2\Gamma=\Gamma^+\setminus{}_2\Gamma^+
\]

donc, utilisant e)

\[
\deleted{\unsure{\Gamma^+}}\quad\Gamma^+=\text{ss-gpe de }\Gamma\text{ engendré par }\Gamma\setminus{}_2\Gamma\tag*{$(*)$}
\]

ce qui montre donc que $\Gamma^+$ est déterminé
de façon unique. D'ailleurs, $\Gamma^-$ est un torseur
à gauche sous $\Gamma^+$, et la multiplication sur
$\Gamma=\Gamma^+\amalg\Gamma^-$ se déduit de la structure de
groupe sur $\Gamma^+$ et de la structure de $\Gamma^+$ tors.
à g. sur $\Gamma^-$, grâce à

1°) La multiplication de $\Gamma$ dans $\Gamma^+\times\Gamma^+$ est celle a $\Gamma^+$\\ 
2°) \rule{4em}{.25pt} $\Gamma^-\times\Gamma^-$ est l'application Cte $\longrightarrow1\in\Gamma^+$\\ 
3°) \rule{3em}{.25pt} $\Gamma^+\times\Gamma^-$ est induite par l'opération de $\Gamma^+$ sur $\Gamma^-$\\ 
4°) \rule{3em}{.25pt} $\Gamma^-\times\Gamma^+$ est donnée par $(\sigma,\unsure{v|u})\longrightarrow\sigma v=(-v)\cdot\sigma$

% archive sheet number: 38
